\chapter{Unbounded derived tensor and internal Hom}
\label{ch:unbounded}

Throughout the chapter \(R\) is a commutative ring. Complexes are
cochain complexes, with \(X[s]^i=X^{i+s}\) and
\(d_{X[s]}=(-1)^s d_X\). Modules are placed in degree zero. We use the
construction of bounded-below injective replacements and the homotopy
calculus of \cite[R01, R03, R04]{ha-core}. In particular, no assumption
of Noetherianity is needed in this chapter.

\section{Complexes and free cells}

\begin{definition}[Tensor, internal Hom, and good models]\label{res:U01}
The tensor complex is totalized by direct sums, with
\(d(x\otimes y)=d x\otimes y+(-1)^{|x|}x\otimes d y\).
The internal Hom complex is
\[
 \underline{\mathrm{Hom}}_R(X,Y)^q
   =\prod_i\mathrm{Hom}_R(X^i,Y^{i+q}),\qquad
 \partial f=d_Y f-(-1)^q f d_X .
\]
A complex \(P\) is K-projective when
\(\underline{\mathrm{Hom}}_R(P,A)\) is acyclic for every acyclic \(A\).
A complex \(F\) is K-flat when \(A\otimes_R F\) is acyclic for every
acyclic \(A\). A complex \(I\) is K-injective when
\(\underline{\mathrm{Hom}}_R(A,I)\) is acyclic for every acyclic \(A\).
\end{definition}

Expanding the two differentials proves that their squares are zero.
For the Hom differential, the two terms containing \(d_Y f d_X\)
cancel. Cycles in degree zero are chain maps, and degree-zero
boundaries are null-homotopic maps. The vanishing condition in all
degrees is thus the usual orthogonality to all shifts of acyclic
complexes. Direct sums in the tensor and products in the Hom are
essential when both arguments are unbounded. For example the
degree-\(q\) tensor--Hom adjunction is the correspondence
\[
 u(x\otimes y)\ \longleftrightarrow\ [x\longmapsto(y\longmapsto
 u(x\otimes y))].
\]
For \(u\) of degree \(q\), each inner map has degree \(q+|x|\).
Expanding \(\partial u\) gives the three terms acting on the target,
on \(x\), and on \(y\), with coefficients \(1,-(-1)^q\),
and \(-(-1)^{q+|x|}\), respectively. This is exactly the Hom
differential on the curried expression. The constructions and the
adjunction carry chain homotopies to chain homotopies.

\begin{theorem}[A semifree replacement without bounds]\label{res:U02}
Every complex \(X\) admits a degreewise surjective quasi-isomorphism
\(F\to X\) such that
\(F\) has an increasing filtration by subcomplexes
\(0\subset F_0\subset F_1\subset\cdots\) with union \(F\).
Each inclusion is split on the underlying graded modules and the
successive quotients are sums of shifts of \(R\).
The complex \(F\) is both K-projective and K-flat.
\end{theorem}
\begin{proof}
View \(R\) as a one-object dg category in degree zero.
Its representable module is the complex \(R[0]\).
Theorems~\ref{res:G06} and \ref{res:G07} give the
graded-split filtration, quasi-isomorphism, and K-projectivity.

Tensor the filtration with \(A\). Splitting on graded modules makes
each successive sequence exact as complexes. Each quotient is a
direct sum of shifts of \(A\), hence acyclic. The finite-stage tensor
complexes are therefore acyclic by induction and their filtered
union remains so, since filtered colimits are exact for modules.
Thus \(F\) is K-flat.
\end{proof}

\begin{proposition}[Invariance and bounded tests]\label{res:U03}
The functors \(P\otimes_R-\), \(\underline{\mathrm{Hom}}_R(P,-)\),
and \(\underline{\mathrm{Hom}}_R(-,I)\) preserve quasi-isomorphisms
when \(P\) is K-flat, \(P\) is K-projective, and \(I\) is
K-injective, respectively. A bounded-above complex of flat modules
is K-flat. Every K-projective complex is K-flat.
\end{proposition}
\begin{proof}
The cone of a quasi-isomorphism is acyclic. Tensor takes a cone to
the cone of the tensor map, up to the canonical shift sign.
In the component involving \(x\in P^i\) and the shifted input the
comparison is \((-1)^i\), as is seen by expanding the differential
in Definition~\ref{res:U01}. Internal Hom in either variable
gives the corresponding cone up to shift, since the cone uses
a graded finite direct sum and its differential contains the
original map as off-diagonal entry. Thus the three acyclicity
conditions imply the asserted invariance.

For a bounded-above complex of flats, first cut off the degrees
below \(m\) by the subcomplex consisting of its terms in degrees
at least \(m\). This is a bounded complex filtered by degree tails,
whose one-term quotients are shifts of flat modules. Tensoring with
an acyclic complex preserves those graded-split exact sequences,
and each one-term quotient gives an acyclic tensor by flatness.
As \(m\) decreases the union is the original complex, so exactness
of filtered colimits proves K-flatness. Finally, apply
Theorem~\ref{res:U02} to a K-projective complex. The quasi-isomorphism
from its semifree replacement is between K-projectives, hence is a
homotopy equivalence by \cite[R04]{ha-core}. Tensor with an acyclic
complex preserves that homotopy equivalence. 
\end{proof}

\section{Injective replacements from a tower}

\begin{lemma}[Strict extension to an injective model]\label{res:U04}
Let \(i:X\hookrightarrow J\) be a termwise monic quasi-isomorphism.
If \(I\) is a K-injective complex with injective terms, every
chain map \(f:X\to I\) extends to a chain map \(J\to I\).
\end{lemma}
\begin{proof}
Extend the components first, using injectivity, to obtain a graded
map \(u:J\to I\) of degree zero. Since \(u i=f\), the degree-one
map \(\partial u\) vanishes on \(X\) and factors through \(J/X\)
as a cycle in \(\underline{\mathrm{Hom}}_R(J/X,I)\).
The quotient \(J/X\) is acyclic by the long exact sequence.
K-injectivity makes this cycle the differential of a degree-zero
map \(v:J/X\to I\). Subtract \(v\) composed with the quotient
map from \(u\). The result extends \(f\) and is closed.
\end{proof}

\begin{lemma}[Countable towers of complexes]\label{res:U05}
For a tower of complexes \(A_0\leftarrow A_1\leftarrow\cdots\)
whose transition maps are degreewise surjective, there is a short
exact sequence of complexes
\[
 0\longrightarrow\varprojlim_n A_n\longrightarrow
 \prod_{n\geq0} A_n
 \xrightarrow{\ 1-s\ }\prod_{n\geq0} A_n
 \longrightarrow0,
\qquad
 s((a_n)_n)=(t_n(a_{n+1}))_n .
\]
A degreewise surjective tower of acyclic complexes has acyclic
limit. A map between degreewise surjective towers that is a
quasi-isomorphism at each stage gives a quasi-isomorphism of
limits. If all \(A_n\) are K-injective and transitions split
in each degree, then the limit is K-injective.
\end{lemma}
\begin{proof}
In each degree the kernel is the compatible families.
To solve \(a_n-t_n(a_{n+1})=b_n\), choose \(a_0\) and then
lift \(a_n-b_n\) recursively. Products of modules are exact,
since surjectivity can be tested in each component. They
therefore commute with taking cohomology of complexes.
The first two consequences follow from the long cohomology
sequences of the displayed short exact sequence.
For the last one, if \(C\) is acyclic pass to the tower
\(\underline{\mathrm{Hom}}_R(C,A_n)\). Its terms are acyclic,
its transition maps are surjective by the degreewise splittings,
and its limit is \(\underline{\mathrm{Hom}}_R(C,\varprojlim A_n)\)
since limits commute with products and with Hom in the target.
Apply the first consequence.
\end{proof}

\begin{theorem}[Unbounded K-injective replacement over a ring]\label{res:U06}
Every complex \(X\) of \(R\)-modules admits a termwise monic
quasi-isomorphism to a K-injective complex with injective terms.
\end{theorem}
\begin{proof}
Put \(T_n=\tau_{\geq -n}X\), where the cutoff term in degree
\(-n\) is \(X^{-n}/d X^{-n-1}\). The maps
\(T_{n+1}\to T_n\) are surjective in each degree, and for
each fixed degree the tower stabilizes at the corresponding
term of \(X\). Thus \(\varprojlim T_n=X\).

We construct compatible termwise monic quasi-isomorphisms
\(T_n\to I_n\) to bounded-below complexes of injectives
whose transitions are split epimorphisms in each degree.
Start with a bounded-below replacement of \(T_0\) by
\cite[R01]{ha-core}. Given \(I_n\), choose a bounded-below
replacement \(T_{n+1}\hookrightarrow J\).
The composite \(T_{n+1}\to T_n\to I_n\) extends to \(g:J\to I_n\)
by Lemma~\ref{res:U04} and bounded-below K-injectivity
\cite[R03]{ha-core}. Put \(C=\mathrm{Cone}(1_{I_n})[-1]\).
With degree-\(i\) term \(I_n^i\oplus I_n^{i-1}\), its differential is
\((a,b)\mapsto(d a,-a-d b)\). Projection onto the first entry is a
chain map \(C\to I_n\) and is surjective with a graded section.
The homotopy \((a,b)\mapsto(-b,0)\) contracts \(C\).
Use \(I_{n+1}=J\oplus C\), map \(T_{n+1}\) into the first summand,
and take transition \(g+\mathrm{pr}_1\).
The map from \(T_{n+1}\) is monic and a quasi-isomorphism,
and commutativity with the transition to \(I_n\) is strict.

Let \(I=\varprojlim I_n\).
Lemma~\ref{res:U05} proves that the natural map \(X\to I\)
is a quasi-isomorphism and that \(I\) is K-injective.
It is monic termwise by compatibility with all the monomorphisms
\(T_n\hookrightarrow I_n\) and \(\varprojlim T_n=X\).
For each degree \(i\), choose splittings of the transition maps
\(I_{n+1}^i\to I_n^i\). Then \(I^i\) is isomorphic as a module to
\(I_0^i\) times the product of their kernels. Each kernel is a direct
summand of an injective module. A product of injectives is injective,
since one can extend a map along a monomorphism in each coordinate.
This proves the final assertion.
\end{proof}

\section{The closed derived structure}

\begin{theorem}[Unbounded derived tensor]\label{res:U07}
There is a symmetric monoidal product
\(\otimes_R^{\mathbf L}:D(R)\times D(R)\to D(R)\) with unit \(R[0]\).
It is computed by replacing both arguments with the models in
Theorem~\ref{res:U02}. If either argument is already K-flat, one
may use the ordinary tensor product without replacing the other.
\end{theorem}
\begin{proof}
By \cite[R04]{ha-core}, K-projective models lift all derived
morphisms uniquely as homotopy classes of chain maps.
Tensoring those classes produces a bifunctor taking values in
\(D(R)\), and the lifted equivalences between models provide its
natural comparison isomorphisms. If \(P,Q\) are K-flat, then
\(P\otimes Q\) is K-flat: associate its tensor with an acyclic
complex and use acyclicity successively. The ordinary chain
associativity, unit and symmetry
\(x\otimes y\mapsto(-1)^{|x||y|}y\otimes x\)
therefore descend, and satisfy the coherence identities because
they already do so on the chain models. Alternatively, after
replacing one of two K-flat complexes by a K-projective model,
the ordinary comparison remains a quasi-isomorphism by the
K-flatness of the other, by Proposition~\ref{res:U03}.
This proves the step needed to iterate the bifunctor independently
of choices. For a K-flat factor tensoring the replacement map in
the other factor is a quasi-isomorphism, again by that proposition.
\end{proof}

\begin{theorem}[Internal derived Hom and adjunction]\label{res:U08}
Write \(\mathbf R\mathrm{Hom}_R(X,Y)\in D(R)\) for internal
derived Hom. It can be computed by
\(\underline{\mathrm{Hom}}_R(P_X,Y)\) with \(P_X\) K-projective,
or by \(\underline{\mathrm{Hom}}_R(X,I_Y)\) with \(I_Y\)
K-injective. Both agree naturally, and
\[
 \mathrm{Hom}_{D(R)}(Z\otimes_R^{\mathbf L}X,Y)
  \cong \mathrm{Hom}_{D(R)}(Z,\mathbf R\mathrm{Hom}_R(X,Y)).
\]
In particular
\(
 H^q\mathbf R\mathrm{Hom}_R(X,Y)
   \cong\mathrm{Hom}_{D(R)}(X,Y[q]).
\)
\end{theorem}
\begin{proof}
Take the two replacements. The expression
\(\underline{\mathrm{Hom}}_R(P_X,I_Y)\) defines a bifunctor
on derived morphisms by \cite[R04]{ha-core} and passage of Hom
through chain homotopy classes. The two single replacement
expressions compare to this one by quasi-isomorphisms by
Proposition~\ref{res:U03}. If \(P\) is K-flat and \(I\) is
K-injective, then \(\underline{\mathrm{Hom}}_R(P,I)\) is
K-injective. Indeed its Hom from an acyclic \(A\) identifies
as a complex with the Hom from \(A\otimes P\) to \(I\), by the
adjunction of Definition~\ref{res:U01}. That complex is acyclic.
For the adjunction take \(P_X,I_Y\) and use the ordinary
chain tensor--Hom adjunction in the homotopy category.
On the left, \(Z\otimes P_X\) computes derived tensor by
Theorem~\ref{res:U07}, and on the right the Hom target is
K-injective by the preceding argument. Both morphism sets
are therefore derived morphism sets by \cite[R04]{ha-core}.
The last assertion follows from the cycle and boundary description
in Definition~\ref{res:U01}, also with the shift of \(Y\).
\end{proof}

\begin{proposition}[Derived extension of scalars]\label{res:U09}
For a homomorphism of commutative rings \(R\to S\), put
\(L(X)=S\otimes_R^{\mathbf L}X\in D(S)\).
Then \(L\) is left adjoint to restriction of scalars, and has the
natural symmetric monoidal comparison
\[
 L(X\otimes_R^{\mathbf L}Y)\ \cong\
 L(X)\otimes_S^{\mathbf L} L(Y).
\]
For successive changes of commutative rings \(R\to S\to T\), the two left extensions
compose to \(T\otimes_R^{\mathbf L}(-)\).
\end{proposition}
\begin{proof}
If \(P\) is K-flat over \(R\), \(S\otimes_R P\) is K-flat over \(S\)
by associativity tested against acyclic \(S\)-complexes.
Similarly extension of a K-projective is K-projective, using the
underived Hom adjunction against an acyclic \(S\)-complex and
exactness of restriction. On a K-projective source the chain Hom
adjunction therefore computes the derived morphisms on both sides.
Choose semifree \(P_X,P_Y\) for the monoidal formula. Both tensors
are computed by the displayed model comparisons, and their
identification is the associativity isomorphism for ordinary
change of rings. It commutes with symmetry, multiplication and the
unit on the chain models, since \(S\) is concentrated in degree zero.
The same base-changed models give transitivity.
\end{proof}
